\documentclass{jarticle}

\title{簡易木製スピニングリール制作について}

\begin{document}


\maketitle

\section{初めに}
この文章を通読する必要はない.
写真だけ,　設計図だけ見て作り始めるのも可.
試行錯誤もまた工作の楽しみの一つである.

\subsection{}
\begin{enumerate}
\item 制作過程を楽しむ(楽しめる)ことがわりあい大切.
\item 市販品の性能を求めては落胆する. ロマンを重視しよう.
\end{enumerate}

※余談になるが, ある程度性能のあるリールをできるだけ安く手に入れたいという方向に効率を求める場合は
釣り人の知人を作ることを進める.
趣味人というものは, まだまだ使える道具があっても新しいものに手を出してしまいがちで,
性能も状態も十分な道具を余らせがちなものである.


設計するにあたって「スピニングリール」を
	回転するアームと, 竿と平行に固定された糸巻を持ち,
	次の機構を持っているものととらえている.
	\begin{enumerate}
	\item
		糸を回転するアームに引っ掛けて, 糸巻に巻きつけることによって巻き取る.
	\item 
		アームから糸を外すことにより, 糸の自由放出(とそれを利用した仕掛けの投擲)が可能.
	\end{enumerate}
	制作難度の調整結果としてペールやオシレーション機構,ドラグ機構等は省略している.
	\begin{enumerate}
	\item
		ペールの省略:糸の引っ掛けが手動なのでもたつく→がんばれ
	\item 
		オシレーション機構の省略:糸絡みの可能性の増大→そんなにヘビーに使うもんではない
	\item 
		ドラグ機構の省略:ロマンが減る→ロマン成分は自作に求める
	\end{enumerate}
\subsection{}
目標
\begin{enumerate}
\item 木製
\item 実用可のスピニングリール
\item 簡単に作れる
\end{enumerate}
目標の達成度
\begin{enumerate}
\item 木製:竿への取り付け部を除いて全木製
\item 実用可:一応, 投げも可. 小魚程度なら問題なく揚げれるし竿でうまく寄せてタモですくえばそこそこのサイズでも対応できなくもない.
\item 簡単に作れる:自分的には達成したが, 最終的には作った人が判断してほしい.
\end{enumerate}
目標の設定動機
\begin{enumerate}
\item 木製:手段と目的の反転. 
\item 実用可のスピニングリール
\item 簡単に作れる
\end{enumerate}


\section{制作}
\subsection{材料}
\begin{enumerate}
\item 100mm x 100mm x 6mm MDF板, 4枚（DAISOで6枚組で売ってる）
\item 丸箸 (コンビニ等で弁当とかに付けてくれるやつ)
\item 竹串
\item 木工ボンド
\item 瞬間接着剤
\end{enumerate}
2～5も家のどっかになければ100均でも入手可.
2，5はそれっぽい棒で用意したドリル刃の径と適合してれば何でも代用可.

実際に遣うためには以下も必要
\begin{enumerate}
\item 巻くための釣り糸(ナイロン糸で十分)
\item 竿取付用金具(適当な大きさのステー等で代用可, 何なら竿も自作してねじ止めしても良し)
\end{enumerate}

\subsection{道具}
\begin{enumerate}
\item のこぎり
\item 平彫刻刀あるいは小さなノミ（C2部材の溝切り落としに使うだけ）
\item ラチェットドライバーとドリル刃
	(6mmと2.8mm, 丸箸と竹串の直径に大体あってるもの.
	多少いい加減でも良いが太い方は緩く細い方はきつめの方が都合良い.)
\item コンパスと定規, 筆記具
\item (木工)やすり
\end{enumerate}

あると便利な道具
\begin{enumerate}
\item マスキングテープ（仮固定用）
\item クランプ（固定, 圧着用）
\item 丸やすり（穴径の調整）
\item 垂直穴あけ用のジグ
\item ホビーカンナ(円に削るときとか早くなる)
\item 止型定規（線引きに何かと便利）
\end{enumerate}

\subsection{作り方}
最重要部分はAパーツBパーツの2つの歯車部分で, 特に歯の部分の精度が重要なのでそこだけは気を付けること.
逆を言えばそれ以外の部分は多少適当でも動く.

\subsubsection{部材の切り出し}


円はのこぎりで大雑把に丸くした後, やすりなどで丸くする.
実を言えば機械機構としての円の作成を要求する部位はない(丸棒の丸とドリル穴の丸を使っている)のでそこまで頑張って丸くしなくても良い.
A2B1D1が丸いのは他の部材との干渉を避けるためであり, D2D3が丸いのは糸の出をスムーズにする為であるので
これらをクリアする程度の円でよい.


\subsubsection{組み立て}
\begin{enumerate}
\item
	設計図に従ってABCD各パーツを組む.
	\begin{itemize}
	\item 基本木工ボンドで圧着
	\item AパーツBパーツの歯の根元に最後に瞬間接着剤をたらして強度を上げる
	\end{itemize}
\item
	Dパーツに糸の端を固定する.
\item
	AパーツにDパーツを挿し, Aパーツが回転できるようにしながらDパーツをCパーツに固定する.
	Dパーツは回転してはいけない.
\item
	適当な長さに切った丸箸をBパーツとCパーツに通し,
	回転可能な様に箸の両端を止める.
\end{enumerate}
各パーツどうしの固定は接着しない木釘(というか竹串釘)
を用いるなど取り外し可能な様にできればその方が使い勝手が良い.


\section{バリエーション}
\subsection{木製縛りの解除}

\end{document}

